% Problem Statement and Current-Law Gap. Three connected movements plus one
% source table. Laws are cited by key so each resolves in the bibliography; the
% dense crosswalk is reserved for the prior-law section.

The problem this bill addresses, the body of existing law that bears on it, and
the reason that body falls short are set out below in three movements. The
sources synthesised here are listed once in Table~\ref{tab:problem-sources}; the
exhaustive statutory crosswalk is deferred to
Section~\ref{sec:prior_law} so that the argument here stays readable.

\clearpage

\subsection{The Problem}

Physical AI oncology trials now place autonomous and semi-autonomous robots in
direct physical contact with patients. The contact takes several forms:
surgical humanoids that cut, suture, and manipulate tissue; therapeutic
positioning systems that hold and move a patient during treatment; diagnostic
platforms that place needles for biopsy or delivery; and rehabilitative and
companion systems that support, monitor, and assist. In each form, the behavior
of the robot at the moment it touches the patient is determined by software, and
that software is increasingly machine-generated rather than written line by line
by a human engineer. No current United States law requires the verification of
that software to clear before the software is generated or executed. A sponsor
may generate robot-patient interaction code, run it, and document the result
afterward, in any order it chooses, because the order is not regulated. H.R. 9510 closes this safety gap.

Embodiment is what raises the stakes from a software-quality concern to a
patient-safety mandate. In a disembodied analytics pipeline, an error produces a
wrong number that a human can catch downstream. In an embodied surgical system,
the same class of error produces a wrong motion that a human cannot catch in
time, because one body concentrates all of the error potential into a single
moving instrument acting in real time. The author's first VVUQ study
\cite{kawchak2026vvuq01} develops this argument and demonstrates its operational
consequence, that the decision to permit execution must precede execution rather
than follow it.

A second structural feature of the field deepens the gap. A Physical AI system
used in a trial occupies an ambiguous regulatory position: it is at once a
medical product, when it administers or assists treatment, and a research tool,
when it generates the evidence the trial is designed to collect. The national
platform's regulatory analysis describes this dual classification and the
uncertainty it creates about which authority, and which pathway, governs the
robot's control code \cite{kawchak2026national}. Verification before generation
is attractive precisely because it is agnostic to that classification: whatever
the robot is called, its patient-facing code must clear the gate before it runs.

\subsection{Current-Law Inventory and the Gap}

The existing framework is substantial, and the bill builds on it rather than
displacing it; but each component stops short of a verification-before-generation
mandate for robot-patient interaction code. The Federal Food, Drug, and Cosmetic
Act \cite{usc-fdca} and the medical-device regulations \cite{cfr-devices}
establish device classification and premarket review, but they govern the device
as a marketed article, not the order in which its trial-conduct code is verified
and generated. The investigational pathways for new drugs and for the protection
of human subjects \cite{cfr-part312, cfr-part50}, together with the author's
Physical AI adaptations of those pathways and of good clinical practice
\cite{kawchak2026cfr312, kawchak2026cfr50, kawchak2026iche6r3}, reach the conduct
of the trial and the documentation of robotic systems, but they do not condition
code generation on a cleared verification record. 

The Medicare standard of
reasonableness and medical necessity \cite{usc-1395y} governs whether a procedure
is covered, not whether the robot's code was verified before it ran. The Public
Health Service Act \cite{usc-phsa} supplies public-health and biologics
authority that sits beside, but does not reach, the sequencing question. The
HIPAA Privacy and Security Rules \cite{cfr-hipaa} and the electronic-records rule
\cite{cfr-part11} protect the integrity and confidentiality of the records a
verification gate produces, which is necessary for an auditable gate but is not
itself a gate.

The federal nondiscrimination guarantees complete the inventory, and a citation
correction belongs here. The guarantees against bias on protected
characteristics are codified at title 42 of the United States Code, not title 26:
Title VI of the Civil Rights Act appears at 42 U.S.C. § 2000d
\cite{usc-titlevi}, and the Americans with Disabilities Act at 42 U.S.C. § 12101
\cite{usc-ada}. An informal reference to 26 U.S.C., which is the Internal Revenue
Code, for these provisions was a mislabel; this bill cites the correct titles
throughout. These guarantees require that a robot-patient algorithm and its
training data minimize the risk of bias and respect accessibility, but they do
not require that the algorithm be verified before it is generated. The national
platform's government-framework and regulatory-landscape analyses supply the
three-branch authority map and the current posture of the Food and Drug
Administration toward artificial intelligence on which this inventory draws
\cite{kawchak2026national}. The pattern across the whole framework is consistent:
the substantive content a gate would enforce already exists in scattered form,
but the rule that makes verification a precondition of generation does not.

\subsection{The State Patchwork and the Case for Uniform Federal Action}

Several states have begun to regulate artificial intelligence in health
decisions, but each acts at the utilization-review or claims-handling layer, not
at the robot-patient code-verification layer. New York would require insurers to
submit AI algorithms and training data for certification of minimized bias and
adherence to evidence-based guidelines \cite{ny-a9149}. Texas would require an
annual artificial-intelligence compliance statement, with the underlying
algorithms produced only on a finding of noncompliance \cite{tx-sb1822}.
California keeps medical-necessity determinations with a licensed physician and
bars reliance on generalized rather than individual clinical data
\cite{ca-sb1120}. Florida bars artificial intelligence as the sole basis for a
claim denial and requires independent human review \cite{fl-hb527}. These are
sound measures, and the bill is designed to harmonise with them; but every one
of them addresses how a payer decides a claim, not how a robot's code is cleared
before it acts on a patient in a trial.

The consequence is that the most safety-critical software in the field, the code
that governs a robot's physical contact with a patient during trial conduct, is
the software least addressed by current law, and is addressed by no two
jurisdictions in the same way. A patchwork at the claims layer cannot supply a
uniform rule at the conduct layer, and a sponsor operating across states cannot
build one robot-patient code base against four divergent regimes. A single
federal verification-before-generation standard resolves this directly and gives
the states a stable upstream rule to build on. The detailed mapping of each
federal and state authority to the clause of this bill that extends it is carried
in the crosswalk of Section~\ref{sec:prior_law}.

\begin{table}[ht]
\centering
\begin{tabularx}{\textwidth}{L{7.5cm} Y}
\toprule
\textbf{Source file (abbreviated leaf)} & \textbf{What the section takes from it} \\
\midrule
\multicolumn{2}{l}{\textbf{papers/VVUQ-01/final-paper/sections/}} \\
introduction.tex & Embodiment raises the stakes; the verification-first gap \\
\multicolumn{2}{l}{\textbf{physical-ai .../final\_paper/sections/}} \\
regulatory\_landscape.tex & Medical-product versus research-tool classification; gaps and pathways \\
gov\_framework.tex & Three-branch authority map; current FDA artificial-intelligence framework \\
\bottomrule
\end{tabularx}
\caption{Sources for the problem statement and current-law gap.}
\label{tab:problem-sources}
\end{table}
